% Working draft. Venue not chosen: neutral article layout; swap in the venue style file when one is set.
% Every number here is quoted from docs/growing_bandits/{RESULTS,DEPLOYMENT_RESULTS}.md and the tables named in
% docs/paper/PAPER_PLAN.md (claim ledger, section 1). Do not add a number that is not in the ledger.
\documentclass[11pt]{article}
\usepackage[margin=1in]{geometry}
\usepackage{amsmath,amssymb,amsthm}
\usepackage{booktabs}
\usepackage{graphicx}
\usepackage[numbers,sort&compress]{natbib}
\usepackage{hyperref}
\usepackage{xcolor}
\newcommand{\todo}[1]{\textcolor{red}{[TODO: #1]}}
\newcommand{\Kstar}{K^{\star}}

\title{When Is Prompt Search Worth It?\\[2pt]\large Search Breadth, Reservoir Heterogeneity, and Flat Real Prompt Pools}
\author{Sanjay Mukhyala \and Ian Waudby-Smith}
\date{Draft --- \today}

\begin{document}
\maketitle

\begin{abstract}
Prompt optimization for LLM agents repeatedly faces one choice: evaluate a new candidate prompt, or evaluate an
existing one again. We study this choice as best-arm identification over a growing reservoir of candidates and ask
when it matters. In simulation, the sequential decision collapses to a single quantity, the number of candidates
evaluated: holding that number fixed, the timing of search is worth nothing, and the optimal number is interior and
grows sublinearly with the evaluation budget. A policy learned from 87{,}148 oracle-labelled states reduces to a
three-parameter rule, and a two-constant recruit-then-refine schedule captures nearly all of the attainable value.
How much choosing the number of candidates matters is set by the reservoir: it rises steeply with the true spread of
candidate quality, and at equal spread it is dominated by a rare upper tail. We then measure that spread directly
for three real prompt families on web-agent tasks --- structured stylistic prompts, deliberately diverse free-form
prompts, and procedural prompts derived from product documentation --- using replicate episodes and nonparametric
deconvolution to separate execution noise from true differences. All three pools are flatter than every simulated
environment, none shows heterogeneity that replicates across tasks, and search breadth barely matters on any of
them. For the agent and applications we measured, the bottleneck in adaptive prompt search was not the search
algorithm but a candidate pool whose members behave differently.
\end{abstract}

% =====================================================================================================
\section{Introduction}
\label{sec:intro}

Choosing a system prompt for an LLM agent is an expensive evaluation problem. Each trial is a full agent episode
whose outcome is a noisy binary success, and the set of candidates is never fixed: a designer, or an optimizer,
can always write another prompt. Methods for automatic prompt optimization differ in how they propose
candidates~\citep{zhou2023ape,yang2024opro,fernando2023promptbreeder,khattab2024dspy,agrawal2025gepa}, but every
one of them must decide, at each step, whether to spend the next evaluation on a new candidate or on a candidate
it already holds. This paper asks when that decision matters.

We formalize the problem as best-arm identification over a growing reservoir. A SEARCH action draws a new prompt
from an effectively unlimited distribution; a REFINE action evaluates a held prompt again. At the end of the budget
the procedure names one prompt and is scored by its true success rate. Simple regret then splits exactly into a
\emph{discovery} term --- the best prompt was never drawn --- and a \emph{selection} term --- it was drawn but not
recognized. Searching more lowers the first and raises the second.

Our study has two halves. The first asks \emph{how} to search, in a simulator where candidate quality is known.
Using an oracle that labels each state with the regret-minimizing action by paired replay, we find that the
decision is driven by how good the held candidates are relative to a fresh draw, not by statistical evidence about
the leader; that, holding the final number of candidates $K$ fixed, the timing of search is worth nothing; that the
optimal $K$ is interior and grows sublinearly with the budget; and that a learned policy collapses onto a
three-parameter rule which a simple budget-scaled schedule nearly matches. The value of all of this, however,
depends on the reservoir: it rises steeply with the spread of true candidate quality and, at equal spread, with
the weight of a rare excellent tail.

The second half asks \emph{whether} real prompt pools have that spread. We measured three deliberately different
families on web-agent tasks: structured prompts varying planning, verification and expertise style; free-form
prompts written to be maximally diverse; and procedural prompts derived from the application's user manual, scored
under a per-episode step budget designed to let knowledge show through as efficiency. Replicate episodes identify
execution noise, and a nonparametric deconvolution recovers the distribution of true prompt success rates. All
three pools are flatter than every simulated environment; none shows heterogeneity that replicates across halves of
the task set; and on every pool, every sensible choice of $K$ is within a few thousandths of the best, with only
``open a new prompt every episode'' measurably worse. A pre-registered test of the simulated policy ordering on a
heterogeneous real pool could not be run, because no pool met its precondition.

\paragraph{Contributions.}
\begin{enumerate}
  \item We show that the search-versus-refine decision in prompt search reduces to choosing search breadth $K$,
        characterize how the optimal breadth scales with budget, and show that mis-sizing breadth dominates any
        difference between policies (\S\ref{sec:breadth}).
  \item We show that a two-constant recruit-then-refine schedule captures nearly all of the attainable value: a
        learned policy trained on oracle labels reduces to an interpretable three-parameter rule and adds little
        beyond the schedule once breadth is not capped (\S\ref{sec:simple}).
  \item We map reservoir heterogeneity, and its upper tail in particular, to the regret at stake in choosing
        breadth, giving a test that can be run before committing to search (\S\ref{sec:heterogeneity}).
  \item We introduce a measurement protocol for real prompt pools --- replicate-identified noise, deconvolution,
        variance decomposition and split-half reliability --- and find, under pre-registration, that three
        qualitatively different real pools are too flat for search breadth to matter (\S\ref{sec:real}).
\end{enumerate}

The real-data findings are limited to one agent configuration, two applications and pools of 40--50 prompts; we
state these limits where they bind (\S\ref{sec:discussion}).

% =====================================================================================================
\section{Related work}
\label{sec:related}

\paragraph{Prompt optimization.} Automatic prompt optimizers generate and score candidate instructions with an LLM
in the loop~\citep{zhou2023ape,yang2024opro}, evolve populations of prompts~\citep{fernando2023promptbreeder}, or
compile multi-stage programs whose prompts are tuned against a metric~\citep{khattab2024dspy}; reflective methods
propose candidates from execution traces~\citep{agrawal2025gepa}. These methods differ mainly in the proposal
distribution. Our results concern the evaluation side they share: how many candidates to score, and how often.

\paragraph{Bandits for prompt selection.} Prompt selection under a fixed candidate set has been cast as best-arm
identification~\citep{shi2024triple}. We study the setting where the candidate set grows, which is the
infinitely-many-armed bandit~\citep{berry1997infinite,wang2008infinitely,carpentier2015simple}. Classical results
there prescribe recruiting a number of arms that grows with the horizon at a rate set by the reservoir's upper
tail; our simulations measure that envelope directly and find that the tail's mass, not only its exponent, sets it.

\paragraph{Anytime-valid inference.} Betting-based confidence sequences and e-processes give error control that
holds at every stopping time~\citep{waudbysmith2024betting,ramdas2023gametheoretic}, and have been used for
sequential multi-armed testing. We use them for elimination during refinement and test whether their evidence
should drive the decision to search; inside the training support it does not.

\paragraph{Prompt sensitivity.} LLM outputs can be sensitive to small, meaning-preserving changes in prompt
format~\citep{sclar2024quantifying}. Our measurement separates that sensitivity into execution noise, task-specific
interaction and a stable prompt effect, and finds the stable effect small for the agent and tasks we study.

\paragraph{Empirical Bayes deconvolution.} We recover the distribution of true prompt success rates by the
nonparametric maximum-likelihood estimator~\citep{kiefer1956consistency,jiang2009general}, with interval estimates
for variance components from the modified large-sample method~\citep{graybill1980mls}.

% =====================================================================================================
\section{Prompt search over a growing reservoir}
\label{sec:setup}

\paragraph{Model.} A reservoir is a distribution $\nu$ over success probabilities $\mu\in[0,1]$. With a budget of $T$
evaluations, at each step $t$ a procedure either SEARCHes --- draws $\mu_{K+1}\sim\nu$ and evaluates the new prompt
once --- or REFINEs --- evaluates a held prompt $i\le K_t$ again, observing $Y\sim\mathrm{Bernoulli}(\mu_i)$. At $t=T$ it
recommends one held prompt $\hat\imath$. Simple regret is measured against the best of the first $T$ reservoir draws,
$\mu^{\star}_T$, which is attainable and common to all policies under common random numbers:
\begin{equation}
  R_T \;=\; \mu^{\star}_T-\mu_{\hat\imath}
      \;=\; \underbrace{\left(\mu^{\star}_T-\max_{i\le K_T}\mu_i\right)}_{R_{\mathrm{disc}}}
      \;+\; \underbrace{\left(\max_{i\le K_T}\mu_i-\mu_{\hat\imath}\right)}_{R_{\mathrm{sel}}}.
\end{equation}
\todo{Fig.~1: reservoir, SEARCH/REFINE, recruit-then-refine, the two regret terms.}

\paragraph{Policies.} A \emph{recruit-then-refine} policy draws until it holds $K(T)$ prompts and then only refines,
allocating pulls by an anytime-valid elimination rule~\citep{waudbysmith2024betting} and recommending by posterior
mean. The budget-scaled schedule recruits while $K_t<c\,T^{\alpha}$; the fixed-$K$ policy recruits $K$ at once;
always-search draws a new prompt every pull. Learned policies decide SEARCH or REFINE from 71 observable state
features. \todo{Name the level rule here or in \S\ref{sec:simple}; harness details to App.~A.}

\paragraph{What should drive the decision?} To learn and evaluate policies we label 87{,}148 simulated states with
the advantage of SEARCH over REFINE, estimated by replaying each state's future twice under common random numbers.
Features describing the quality of the held candidates predict the label (best single feature AUC 0.706); features
summarizing e-process evidence about the leader do not (best 0.535), and adding all of them to a clock-only model
changes balanced accuracy by $-0.001$ and $+0.006$, against $+0.052$ for quality features. An oracle that sees the
reservoir's upper tail would score AUC 0.836 on the label, while its observable estimate scores 0.572: the decision
is a reservoir-estimation problem, not an inference problem. These results hold inside the training support; out of
support the evidence features recover a little value (App.~B). \todo{Tighten; move feature lists to App.~B.}

% =====================================================================================================
\section{Search breadth is the decision}
\label{sec:breadth}
% Claim P1, P2. Evidence: k_star_envelope.csv, capmatch_contrasts.csv, capp_policies.csv. Fig 2, Fig 3b.
% Not here: Test B / D, the unpaired cap sweep, the 25 learned variants.

\paragraph{Timing does not matter given breadth.} Holding a policy's final number of candidates fixed and
front-loading its recruitment, on the same episodes, the best learned policy ties its control
($+0.0001$ $[-0.0005,+0.0007]$ over 8 environments) and the validation-tuned gradual schedule is worse
($+0.0009$ $[+0.0005,+0.0012]$, worse in 8 of 8). \todo{Fig.~2b; one sentence on why: early recruitment gives every
arm the full refinement budget.}

\paragraph{The optimal breadth is interior and grows with budget.} Pooled $\Kstar$ is 24, 32, 48, 128 and 256 at
$T=50,100,200,500,1000$ (log-log slope about 0.8), with per-environment $\Kstar$ at $T=1000$ spanning 12--400.
Always-search, which maximizes discovery, has regret 0.346 at $T=1000$ against 0.077 for the tuned schedule.
\todo{Fig.~2a U-curves.}

\paragraph{Mis-sized breadth dominates policy differences.} Under a breadth cap of 64, six policies --- including every
learned one --- are identical to four decimals at $T=1000$ (0.0973); lifting the cap lowers regret by 0.020 for every
tuned schedule (pooled; the environment-level interval includes zero above a cap of 128). \todo{Fig.~3b.}

% =====================================================================================================
\section{Simple schedules suffice}
\label{sec:simple}
% Claim P3, N1, N4. Evidence: h1b_*.csv, capp_contrasts.csv, capc_*.csv, surrogate_validity.csv. Fig 3a, Table 1.

\paragraph{The learned policy is a three-parameter rule.} \todo{Registration ladder in two sentences: the learned
policy beats the tuned schedule (0.005), a per-horizon fixed $K$ and an adaptive schedule, then is matched by
``recruit to $T^{0.75}\exp(4(0.5-\text{level}))$ arms, then refine'': $-0.00002$ $[-0.0007,+0.0007]$, $p=0.96$,
30 environments, $T\le 200$.}

\paragraph{With room to work, the edge is small.} \todo{Level rule minus re-tuned schedule $-0.0014$
$[-0.0045,+0.0017]$; learned policies over-recruit to 380--527 arms uncapped, regret 0.14--0.15 vs 0.076.}

\paragraph{The schedule generalizes; learned signals do not.} \todo{Held-out family, uncapped: schedule within 0.001
of the best single $K$ ($-0.0006$ $[-0.0016,+0.0004]$); learned policy $+0.035$; oracle tail exponent worst.}

\paragraph{Offline skill is not a proxy.} \todo{Spearman $+0.005$ $[-0.53,+0.52]$ over 22 variants; the best
deployed variants rank 19th and 22nd offline.}

% =====================================================================================================
\section{Heterogeneity sets the value of search}
\label{sec:heterogeneity}
% Claim P4. Evidence: calibration.csv. Fig 4 (simulation layer).

\todo{Calibration: Beta pools at level 0.6, spread 0.01--0.15, and flat-bulk-plus-rare-great mixtures, $T\le 200$.
Regret range over $K$ at $T=200$: 0.0013, 0.0053, 0.030, 0.080, 0.151 at spread 0.01, 0.02, 0.05, 0.10, 0.15;
always-search loss 0.002 $\to$ 0.188; schedule within 0.013 of the ceiling throughout. At equal spread the tail
decides: spread 0.036 with 1\% of arms at $+0.3$ gives 0.083, against 0.015 for a Beta pool. Define the flat
threshold (spread $\approx 0.019$) and the decision-relevant tier.}

% =====================================================================================================
\section{Real prompt reservoirs are flat}
\label{sec:real}
% Claim P5, S4--S6, N3. Evidence: emp_*, glk30_*, glk_*, stage0_*. Fig 4 (real points), Fig 5, Table 2.
% Not here: secondary task mixes, diagnostics 1--7 detail, latency window, step-cap deviation (App. F).

\paragraph{Measurement.} \todo{Agent, apps, tasks (30 per pool, stratified by difficulty), replicate cells (300 Gmail,
120 GitLab), within-cell noise from replicate pairs, NPMLE deconvolution, ANOVA components with MLS intervals,
split-half reliability, K-grid replay on the deconvolved reservoir. All pre-registered.}

\paragraph{Three pools.} \todo{Grid (50, Gmail), free-form (50, Gmail), procedural (40, GitLab, 30-step budget).}

\paragraph{All three are flat.} \todo{Deconvolved spread 0.035 / 0.001 / 0.041 (0.009 uncapped), all below the
simulated minimum of 0.075; regret range at $T=200$ 0.009 / 0.0001 / 0.008; no detectable heterogeneity (split-half
$r_{\mathrm{SB}}$ 0.07, $p=0.43$; $\le 0$, $p=0.61$); task variance 0.164 vs prompt about 0.001; raw spreads about
twice the true spread. Pools sit below the Beta curve at their own spread: no upper tail.}

\paragraph{A budget exposes efficiency, not knowledge.} \todo{0.036 at 30 steps vs 0.011 uncapped; mean steps vs
rate $\rho=-0.40$; does not replicate across task halves.}

\paragraph{The pre-registered test.} \todo{State the gate and that it did not open; Stage B computed and not read;
directionally always-search $\approx 0.01$ worse and $K=8$ at the minimum.}

% =====================================================================================================
\section{Discussion and limitations}
\label{sec:discussion}

\todo{Practitioner rule: measure deconvolved spread with replicates before searching; flat $\Rightarrow$ small fixed
set; tail $\Rightarrow$ budget-scaled $K$; never always-search. Limits: one agent configuration; two apps; 30 tasks
per pool ($\pm 0.015$--$0.02$ resolution); 40--50 prompts cannot exclude rare excellent prompts below about 2\%
frequency, precisely the regime \S\ref{sec:heterogeneity} shows matters most; replay rather than live runs; synthetic
simulation families; no DP optimum. Future work: prompt sensitivity of weaker or more constrained agents.}

\bibliographystyle{plainnat}
\bibliography{refs}

\appendix
\section{Setup and harness}\todo{}
\section{Oracle labels, commitment horizons, and the e-process ablation}\todo{}
\section{Learned policies and the registration ladder}\todo{}
\section{Breadth caps and generalization}\todo{}
\section{Calibration details}\todo{}
\section{Real-pool protocol, diagnostics, and deviations}\todo{}
\section{Uniform versus adaptive allocation over fixed hand-written prompts}\todo{}
\section{Pre-registration index}\todo{}
\section{Reproduction}\todo{}

\end{document}
